\section*{Chapitre \ref{ch:b1:layered-earth} --- La Terre en couches}

\begin{solution}{exo:b1:layered-earth:1}
Les fractions volumiques sont les cubes des rapports des rayons~: noyau
$(3480/6371)^3 = 0.1630$, soit \qty{16.3}{\percent}~; \omterm{def:b1:layered-earth:layers}{noyau interne}
$(1221.5/6371)^3 = 0.00705$, \qty{0.70}{\percent}.
\end{solution}

\begin{solution}{exo:b1:layered-earth:2}
Au-dessus~: $\Phi = 13.7166^2 - \tfrac43\times7.26466^2 = 188.145 - 70.367 =
\qty{117.78}{km^2.s^{-2}}$, $K = \rho\Phi = \qty{5566.45}{kg.m^{-3}}\times
\qty{117.78e6}{m^2.s^{-2}} = \qty{655.6}{GPa}$. Au-dessous~: $\Phi = 8.06482^2 =
\qty{65.04}{km^2.s^{-2}}$, $K = 9903.49\times\num{65.04e6} = \qty{644.1}{GPa}$.
Le module d'incompressibilité ne change guère (\qty{-1.7}{\percent}). Ce qui
change, c'est le module de cisaillement, $\mu = \rho v_S^2 = \qty{293.8}{GPa}$
au-dessus et $0$ au-dessous (un liquide), et la masse volumique, qui augmente de
\qty{78}{\percent}~: d'où la chute de $v_P$ malgré un $K$ presque égal.
\end{solution}

\begin{solution}{exo:b1:layered-earth:3}
$\sin i_c = 6.2/8.1 = 0.7654$, $i_c = \ang{49.95}$. Distance critique
$2\times30\times\tan\ang{49.95} = 60\times1.1906 = \qty{71.4}{km}$. \omterm{def:b1:layered-earth:head-wave}{Distance de
croisement} $60\sqrt{14.3/1.9} = 60\times2.743 = \qty{164.6}{km}$.
\end{solution}

\begin{solution}{exo:b1:layered-earth:4}
$g = Gm/r^2 = \num{6.674e-11}\times\num{1.94e24}/(\num{3.48e6})^2 =
\qty{10.69}{m.s^{-2}}$, soit \qty{9}{\percent} de plus qu'à la surface. En
descendant jusqu'au noyau, le rayon au carré tombe à $(3480/6371)^2 = 0.298$ de
sa valeur en surface, tandis que la masse attirante ne tombe qu'à
$1.94/5.972 = 0.325$ du total, parce que le noyau contient le tiers de la
masse~: $0.325/0.298 = 1.09$.
\end{solution}

\begin{solution}{exo:b1:layered-earth:5}
Près de la source, $v_1 = 60/10.0 = \qty{6.0}{km.s^{-1}}$. Au loin,
$v_2 = (300 - 200)/(45.8 - 33.3) = 100/12.5 = \qty{8.0}{km.s^{-1}}$~; le temps
d'interception est $t_i = 33.3 - 200/8.0 = \qty{8.3}{s}$. Alors $\sin i_c =
0.75$, $\cos i_c = 0.6614$, et $h = t_i v_1/(2\cos i_c) = 8.3\times6.0/1.3229
= \qty{37.6}{km}$. Vérification~: $x_c = 2\times37.6\times\sqrt{14/2} = \qty{199}{km}$,
où les deux ondes arrivent à $199/6.0 = 199/8.0 + 8.3 = \qty{33.2}{s}$.
\end{solution}

\begin{solution}{exo:b1:layered-earth:6}
Pour $\rho$ uniforme, $I = \tfrac{8\pi}{3}\rho\int_0^R r^4\,\mathrm dr =
\tfrac{8\pi}{15}\rho R^5$ et $M = \tfrac43\pi\rho R^3$, donc
$I/MR^2 = (8/15)/(4/3) = 0.4$. Deux couches ($r_c = \qty{3.48e6}{m}$,
$R = \qty{6.371e6}{m}$)~: $M = \tfrac43\pi\bigl(11000\,r_c^3 + 4450\,(R^3 -
r_c^3)\bigr) = \qty{5.977e24}{kg}$, \qty{0.08}{\percent} au-dessus de la masse
mesurée~; $I = \tfrac{8\pi}{15}\bigl(11000\,r_c^5 + 4450\,(R^5 - r_c^5)\bigr)$
donne $I/MR^2 = 0.3457$. La masse est juste, mais le facteur est trop élevé de
\qty{4.5}{\percent}~: le \omterm{def:b1:layered-earth:layers}{manteau} réel n'est pas uniforme, sa masse volumique
croît avec la profondeur, ce qui déplace de la masse vers l'intérieur.
\end{solution}

\begin{solution}{exo:b1:layered-earth:7}
$\Phi = 12.16^2 - \tfrac43\times6.673^2 = 147.87 - 59.37 =
\qty{88.49}{km^2.s^{-2}}$. Avec la profondeur $z$,
$\mathrm d\rho/\mathrm dz = \rho g/\Phi = 4860\times9.9/\num{88.49e6} =
\qty{5.44e-4}{kg.m^{-4}}$, soit \qty{0.544}{kg.m^{-3}} par kilomètre, ou
\qty{0.54}{g.cm^{-3}} par \qty{1000}{km}. PREM augmente de $5.357 - 4.844 =
0.513$ sur ces \qty{1000}{km}, soit \qty{6}{\percent} de moins~: le gradient
calculé au sommet est appliqué à tout l'intervalle, alors que $\Phi$ croît avec
la profondeur et que le gradient diminue.
\end{solution}

\begin{solution}{exo:b1:layered-earth:8}
$P \approx \bar\rho g z = 3570\times9.9\times\num{670e3} = \qty{23.7}{GPa}$
(PREM~: \qty{23.9}{GPa}). C'est la pression, environ \qty{24}{GPa}, à laquelle la
ringwoodite se décompose en bridgmanite et ferropériclase~: la discontinuité de
\qty{660}{km}.
\end{solution}

\begin{solution}{exo:b1:layered-earth:9}
(a)~Fer liquide~: $10.90\times1.05 = 11.4$ à $10.90\times1.10 =
\qty{12.0}{g.cm^{-3}}$. (b)~$12.764/12.166 = 1.049$~: un saut de
\qty{4.9}{\percent} dans PREM, un peu plus que les \qty{4.5}{\percent} des
études récentes. (c)~La congélation seule ne rendrait le solide que légèrement
plus dense que le liquide~; le saut plus grand signifie que le solide contient
moins d'éléments légers que le liquide~: ils préfèrent le liquide, et la
congélation à la \omterm{def:b1:layered-earth:discontinuity}{limite du noyau interne} en laisse la plupart derrière elle.
\end{solution}

\begin{solution}{exo:b1:layered-earth:10}
Posons $x_* = 2h\tan i_c$. Alors $\sqrt{x_*^2 + 4h^2} = 2h/\cos i_c$ et
$t_r(x_*) = 2h/(v_1\cos i_c)$. L'\omterm{def:b1:layered-earth:head-wave}{onde conique} donne
$t(x_*) = 2h\tan i_c/v_2 + 2h\cos i_c/v_1 = \frac{2h}{v_1}\bigl(\sin^2 i_c/\cos
i_c + \cos i_c\bigr) = 2h/(v_1\cos i_c)$, en utilisant $\sin i_c = v_1/v_2$~: la
même. La pente de l'hyperbole est $t_r'(x) = x/(v_1\sqrt{x^2 + 4h^2})$~; en $x_*$,
elle vaut $\sin i_c/v_1 = 1/v_2$, la pente de la droite de l'\omterm{def:b1:layered-earth:head-wave}{onde conique}. Comme
$t_r'' > 0$, l'hyperbole est convexe et située au-dessus de chacune de ses
tangentes~: pour $x > x_*$, $t_r > t_{\text{conique}}$. À la distance critique, la
réflexion frappe la limite exactement sous $i_c$ et c'est la naissance de l'\omterm{def:b1:layered-earth:head-wave}{onde
conique}~; au-delà, l'\omterm{def:b1:layered-earth:head-wave}{onde conique}, qui gagne en courant dans le milieu rapide,
arrive la première.
\end{solution}

\begin{solution}{exo:b1:layered-earth:11}
$m(r) = \tfrac43\pi\rho r^3$, donc $g = Gm/r^2 = \tfrac43\pi G\rho r$. Alors
$P(r) = \int_r^R \rho\,\tfrac43\pi G\rho s\,\mathrm ds = \tfrac23\pi G\rho^2(R^2
- r^2)$. Au centre, $P = \tfrac23\pi\times\num{6.674e-11}\times5513^2\times
(\num{6.371e6})^2 = \qty{172}{GPa}$, soit \qty{47}{\percent} des
\qty{363.85}{GPa} de PREM~: concentrer la masse vers le centre double à peu près
la pression centrale.
\end{solution}

\begin{solution}{exo:b1:layered-earth:12}
Avec $z$ la profondeur, $\mathrm d\rho/\mathrm dz = \rho g/\Phi$~; si $g/\Phi$ est
une constante $a$, $\mathrm d(\ln\rho)/\mathrm dz = a$ et $\rho = \rho_0
e^{a(z - z_0)}$. Ici $a(z - z_0) = 7.9\times\num{2258.5e3}/\num{87e6} =
0.2051$ et $\rho = 9.90\times e^{0.2051} = 9.90\times1.2276 =
\qty{12.15}{g.cm^{-3}}$, soit \qty{0.2}{\percent} sous les $12.17$ de PREM~: toute
l'augmentation de masse volumique à travers le \omterm{def:b1:layered-earth:layers}{noyau externe} est de
l'autocompression.
\end{solution}

\begin{solution}{pb:b1:layered-earth:1}
\textbf{1.} $V = \tfrac43\pi(\num{6.371e6})^3 = \qty{1.0832e21}{m^3}$.

\textbf{2.} $\tfrac43\pi(r_2^3 - r_1^3)$~: \omterm{def:b1:layered-earth:layers}{croûte}, \omterm{def:b1:layered-earth:layers}{manteau} supérieur et zone de
transition \qty{3.0706e20}{m^3}~; \omterm{def:b1:layered-earth:layers}{manteau} inférieur \qty{5.9961e20}{m^3}~; \omterm{def:b1:layered-earth:layers}{noyau
externe} \qty{1.6890e20}{m^3}~; \omterm{def:b1:layered-earth:layers}{noyau interne} \qty{7.634e18}{m^3}.

\textbf{3.} \qty{28.35}{\percent}, \qty{55.36}{\percent}, \qty{15.59}{\percent}
et \qty{0.70}{\percent}~: le \omterm{def:b1:layered-earth:layers}{manteau} inférieur est l'enveloppe la plus grande, plus
de la moitié du volume de la Terre.

\textbf{4.} $15.59 + 0.70 = \qty{16.3}{\percent}$.

\textbf{5.} \omterm{def:b1:layered-earth:layers}{Croûte}, \omterm{def:b1:layered-earth:layers}{manteau} supérieur et zone de transition $\num{3.0706e20}\times3570 =
\qty{1.0962e24}{kg}$~; \omterm{def:b1:layered-earth:layers}{manteau} inférieur \qty{2.9381e24}{kg}~; \omterm{def:b1:layered-earth:layers}{noyau externe}
\qty{1.8410e24}{kg}~; \omterm{def:b1:layered-earth:layers}{noyau interne} \qty{9.841e22}{kg}.

\textbf{6.} \qty{5.9737e24}{kg}, \qty{0.03}{\percent} au-dessus de la valeur mesurée
\qty{5.972e24}{kg}.

\textbf{7.} $\num{5.9737e24}/\num{1.0832e21} = \qty{5515}{kg.m^{-3}}$.

\textbf{8.} $\num{9.841e22}/\num{5.9737e24} = \qty{1.65}{\percent}$.

\textbf{9.} $(1.0962 + 2.9381)/5.9737 = \qty{67.5}{\percent}$.

\textbf{10.} Le noyau représente \qty{16.3}{\percent} du volume et
\qty{32.5}{\percent} de la masse, le double~: sa masse volumique moyenne,
\qty{10986}{kg.m^{-3}}, est environ $2.5$ fois celle du \omterm{def:b1:layered-earth:layers}{manteau}.

\textbf{11.} Une enveloppe de masse volumique uniforme donne
$I = \tfrac{8\pi}{3}\rho\int_{r_1}^{r_2} r^4\,\mathrm dr = \tfrac{8\pi}{15}\rho
(r_2^5 - r_1^5)$.

\textbf{12.} \omterm{def:b1:layered-earth:layers}{Croûte}, \omterm{def:b1:layered-earth:layers}{manteau} supérieur et zone de transition \qty{2.676e37}{kg.m^2}~; \omterm{def:b1:layered-earth:layers}{manteau} inférieur
\qty{4.525e37}{kg.m^2}~; \omterm{def:b1:layered-earth:layers}{noyau externe} \qty{9.27e36}{kg.m^2}~; \omterm{def:b1:layered-earth:layers}{noyau interne}
\qty{5.87e34}{kg.m^2}.

\textbf{13.} Le \omterm{def:b1:layered-earth:layers}{manteau} inférieur~: il a la moitié de la masse, à de grands rayons.
Le \omterm{def:b1:layered-earth:layers}{noyau interne} contient \qty{1.6}{\percent} de la masse, tout entière près de
l'axe~: sa part est de \qty{0.07}{\percent}.

\textbf{14.} $I = \qty{8.1345e37}{kg.m^2}$ et $I/MR^2 = \num{8.1345e37}/
(\num{5.9737e24}\times(\num{6.371e6})^2) = 0.3355$.

\textbf{15.} \qty{1.4}{\percent} au-dessus de $0.3308$, et bien au-dessous de $0.4$~:
le modèle a la bonne concentration de masse vers le centre, à peu près.

\textbf{16.} Oui. Dans chaque enveloppe, la masse volumique réelle est plus élevée
en bas qu'en haut~; une enveloppe uniforme de même masse met trop de cette masse
près de son rayon externe, où elle compte avec $r^2$, et surestime donc $I$.

\textbf{17.} $g = GM/R^2 = \num{6.674e-11}\times\num{5.9737e24}/
(\num{6.371e6})^2 = \qty{9.822}{m.s^{-2}}$.

\textbf{18.} La masse du noyau est $\num{1.8410e24} + \num{9.841e22} =
\qty{1.9394e24}{kg}$, et $g = G m/r_c^2 = \qty{10.69}{m.s^{-2}}$.

\textbf{19.} $(\num{1.0962e24} + \num{2.9381e24})/(\num{3.0706e20} +
\num{5.9961e20}) = \num{4.0343e24}/\num{9.0667e20} = \qty{4450}{kg.m^{-3}}$.

\textbf{20.} $P \approx 4450\times10\times\num{2.891e6} = \qty{128.6}{GPa}$,
\qty{5}{\percent} au-dessous des \qty{135.75}{GPa} de PREM.

\textbf{21.} $\rho_c = \num{1.9394e24}/\num{1.7653e20} = \qty{10986}{kg.m^{-3}}$,
et $\Delta P = \tfrac23\pi\times\num{6.674e-11}\times10986^2\times
(\num{3.48e6})^2 = \qty{204.3}{GPa}$.

\textbf{22.} $128.6 + 204.3 = \qty{332.9}{GPa}$, \qty{8.5}{\percent} au-dessous des
\qty{363.85}{GPa} de PREM. Les deux estimations ignorent l'augmentation de la
masse volumique avec la profondeur à l'intérieur de chaque région, qui place
davantage de masse plus bas, maintient la pesanteur plus élevée en descendant et
relève $\rho g$, le gradient de pression.

\textbf{23.} La masse totale (question 6), le \omterm{def:b1:comparative-planetology:moi}{facteur de moment d'inertie}
(question 15) et la pesanteur de surface (question 17, contre la moyenne mesurée
\qty{9.82}{m.s^{-2}}) sont vérifiés par des mesures de la Terre~; la pesanteur au
noyau et les pressions ne peuvent être comparées qu'à PREM, lui-même un modèle.

\textbf{24.} Le noyau contient \textbf{\qty[text-series-to-math=true]{32.5}{\percent}}
de la masse de la Terre.
\end{solution}
